\documentclass{article}
\usepackage{amsmath, amssymb, amsthm}
\usepackage{hyperref}
\usepackage{geometry}
\geometry{margin=1in}

\title{Erd\H{o}s Problem \#889}
\date{Last updated: 2025-08-31}
\author{Source: erdosproblems.com}

\begin{document}
\maketitle

\noindent\textbf{Status:} OPEN \quad \textbf{No prize}

\noindent\textbf{Tags:} number theory

\medskip
\noindent\textbf{Problem Statement:}

For $k\geq 0$ and $n\geq 1$ let $v(n,k)$ count the prime factors of $n+k$ which do not divide $n+i$ for $0\leq i<k$. Equivalently, $v(n,k)$ counts the number of prime factors of $n+k$ which are $>k$.

Is it true that\[v_0(n)=\max_{k\geq 0}v(n,k)\to \infty\]as $n\to \infty$?

\bigskip
\noindent\textbf{Remarks:}

A question of Erd\H{o}s and Selfridge \cite{ErSe67}, who could only show that $v_0(n)\geq 2$ for $n\geq 17$. More generally, they conjecture that\[v_l(n)=\max_{k\geq l}v(n,k)\to \infty\]as $n\to \infty$, for every fixed $l$, but could not even prove that $v_1(n)\geq 2$ for all large $n$.

This is problem B27 of Guy's collection \cite{Gu04}.

\bigskip
\noindent\textbf{References:}

[ErSe67] Erd\H{o}s, P. and Selfridge, J. L., Some problems on the prime factors of consecutive integers. Illinois J. Math. (1967), 428--430.
 
 [Gu04] Guy, Richard K., Unsolved problems in number theory. (2004), xviii+437.

\bigskip
\noindent\small{Source: \url{https://www.erdosproblems.com/889}}
\end{document}
